% chktex-file 3
% chktex-file 24
% ============================================================
%  ACL 2025 Long Paper — Two-Column Format
%  "The Illusion of Calibration"
%  Inesh Dheer · Varun Gupta · IIIT Hyderabad
% ============================================================
\documentclass[11pt]{article}

% ── Core ACL packages ──────────────────────────────────────
\usepackage[hyperref]{acl}
\usepackage{times}
\usepackage{latexsym}
\usepackage[T1]{fontenc}
\usepackage[utf8]{inputenc}
\usepackage{microtype}
\usepackage{inconsolata}

% ── Math & science ─────────────────────────────────────────
\usepackage{amsmath}
\usepackage{amssymb}
\usepackage{bm}

% ── Tables & figures ───────────────────────────────────────
\usepackage{booktabs}
\usepackage{multirow}
\usepackage{makecell}
\usepackage{graphicx}
\usepackage{subcaption}
\usepackage{float}

% ── Misc ───────────────────────────────────────────────────
\usepackage{xcolor}
\usepackage{enumitem}
\usepackage{url}

% ── Hyperlinks ─────────────────────────────────────────────
% \usepackage[hyperref]{acl} already loads hyperref; \hypersetup below
% customises colour options only and does not re-load the package.
\hypersetup{colorlinks=true,
            linkcolor=blue!70!black,
            citecolor=blue!70!black,
            urlcolor=blue!70!black}

% ── Layout tuning ─────────────────────────────────────────
\emergencystretch=2em
% \hbadness and \vbadness suppressed for drafting; remove before camera-ready
% to expose genuine overfull \hbox warnings.
%\hbadness=10000
%\vbadness=10000

% ── Custom commands ────────────────────────────────────────
\newcommand{\src}{AIROGS}
\newcommand{\tgt}{Chaksu}
\newcommand{\auroc}{\textsc{auroc}}
\newcommand{\se}{Se}
\newcommand{\spec}{Sp}

% ── ACL anonymization toggle ──────────────────────────────
% ACL final mode is the default; use the review option for blind submission.

% ===========================================================
\title{The Illusion of Calibration: Vulnerabilities of Unsupervised\\
       Test-Time Adaptation in Cross-Ethnic Glaucoma Screening\thanks{\url{https://github.com/iDheer/The-Illusion-of-Calibration}}}

\author{Inesh Dheer \\
  IIIT Hyderabad \\
  \texttt{inesh.dheer@research.iiit.ac.in} \\\And
  Varun Gupta \\
  IIIT Hyderabad \\
  \texttt{varun.gup@students.iiit.ac.in}}

\date{}
% ===========================================================

\begin{document}
\maketitle

% ===========================================================
%  ABSTRACT
% ===========================================================
\begin{abstract}
Prior work documents substantial performance degradation and fairness
disparities when retinal AI models are deployed across demographic and 
ethnic populations~\cite{wang2025,fairDomain2024}. This well-established
vulnerability has motivated the adoption of source-free domain
adaptation (SFDA) as the standard label-free remedy for cross-ethnic
clinical deployment. In this study, we empirically evaluate this assumption
using a modern vision foundation model. We find that a DINOv3-based
classifier trained on the European \src{} dataset achieves an unexpected
$0.852$ \auroc{} zero-shot on Chaksu, an Indian multi-camera fundus dataset,
falling within $0.040$ of a fully supervised oracle ($0.892$) without any 
access to target labels during training. Consequently, no adaptation is 
required. Paradoxically, when we apply standard source-free information-maximisation 
adaptation to further improve performance, the model catastrophically collapses 
to $0.390$ \auroc{} (below chance), with $100\%$ sensitivity
and $1.4\%$ specificity, predicting glaucoma for every patient.
Critically, the SHOT-IM objective is numerically satisfied throughout this
collapse; per-sample entropy and batch-mean prevalence both appear
well-behaved. We term this failure mode the \textit{illusion of
calibration}: a condition where a model satisfies aggregate statistical objectives
while completely losing individual-level diagnostic discrimination. Feature-space
and attention analyses confirm that the adaptation process reorganises
representations around camera-specific artefacts rather than optic-nerve
pathology. Our findings demonstrate that foundation models can inherently generalise across
ethnic and device shifts, and that applying unsupervised adaptation in such 
settings undermines supervised pre-training, constituting a patient
safety risk rather than a deployment solution.
\end{abstract}

% ===========================================================
%  1.  INTRODUCTION
% ===========================================================
\section{Introduction}
\label{sec:intro}

Glaucoma is the second leading cause of irreversible blindness worldwide,
affecting more than 76 million people in 2020 and projected to reach
111.8 million by 2040~\cite{tham2014}. While deep learning models for
automated fundus screening achieve strong within-dataset performance,
deploying them across international demographics introduces a critical
bottleneck: \emph{cross-ethnic domain shift}.

\paragraph{The cross-ethnic generalisation concern.}
It is a widely held assumption that retinal AI models trained on Western 
cohorts will degrade significantly when applied to South Asian or other 
demographically distinct populations. Burlina et al.\ \shortcite{burlina2021} reported a $12.5$ percentage-point
accuracy gap between lighter- and darker-skinned diabetic-retinopathy
patients; FairDomain~\cite{fairDomain2024} documented \auroc{} disparities
of up to $0.085$ between racial subgroups for glaucoma classification. 
More recently, Wang et al.\ \shortcite{wang2025} demonstrated that domain shifts 
in diabetic retinopathy screening induce substantial unfairness, with performance 
gaps disproportionately affecting specific ethnicities. Coyner et al.\ \shortcite{coyner2023} showed that
race-linked signals persist in retinal representations even after vessel-map
abstraction. Rajesh et al.\ \shortcite{rajesh2025} reframed this through the Retinal
Pigment Score, showing that the image-level distance between Western and
Indian retinas reflects a continuous biological gradient rather than a
discrete ethnic boundary. For Indian deployment, the shift is further
compounded by darker retinal pigmentation, non-mydriatic handheld cameras
(Remidio, Forus 3nethra, Bosch), and device-specific artefacts such as
glare, blur, and visible watermarks~\cite{kumar2023}. Figure~\ref{fig:domainshift} illustrates
this compound shift.

\begin{figure*}[t]
\centering
\begin{subfigure}[t]{0.46\textwidth}
  \centering
  \includegraphics[width=\linewidth]{Context/AIROGS.jpg}
  \caption{\src{}}
\end{subfigure}\hfill
\begin{subfigure}[t]{0.46\textwidth}
  \centering
  \includegraphics[width=\linewidth]{Context/Remidio.JPG}
  \caption{Chaksu \textendash{} Remidio}
\end{subfigure}

\vspace{0.6em}
\begin{subfigure}[t]{0.46\textwidth}
  \centering
  \includegraphics[width=\linewidth]{Context/Forus.png}
  \caption{Chaksu \textendash{} Forus}
\end{subfigure}\hfill
\begin{subfigure}[t]{0.46\textwidth}
  \centering
  \includegraphics[width=\linewidth]{Context/bosch.JPG}
  \caption{Chaksu \textendash{} Bosch}
\end{subfigure}
\caption{Compound phenotype and device shift between source and target
  domains. \src{} images are square, well-illuminated tabletop captures.
  Chaksu images come from three handheld devices and exhibit darker
  pigmentation, variable field-of-view, glare, and visible device
  artefacts.}
\label{fig:domainshift}
\end{figure*}

\paragraph{Source-free adaptation as the proposed remedy.}
Because the cross-ethnic failure narrative is widely accepted, source-free domain adaptation (SFDA) and
test-time adaptation (TTA) are routinely proposed as the necessary deployment
remedies. SHOT~\cite{liang2020} established information maximisation and
self-supervised pseudo-labelling as the canonical SFDA objective;
TENT~\cite{wang2021tent}, EATA~\cite{niu2022eata}, and
SAR~\cite{niu2023sar} extended and stabilised it. Their shared
premise (that unlabelled target data should help recover performance degraded by domain shift) is rarely tested directly against modern
self-supervised foundation backbones.

\paragraph{This study.}
Operating under the assumption that the source model would fail on Indian retinas, we designed an experimental pipeline to evaluate whether SFDA could bridge the ethnic gap without requiring expensive target annotations.
Using a DINOv3 ViT-L/16 backbone trained on the European
\src{}~\cite{devente2024} dataset and evaluated on the held-out
Chaksu~\cite{kumar2023} test split (336 images: 51 glaucoma, 285 normal),
we evaluated five conditions:
\begin{enumerate}[leftmargin=1.2em,topsep=2pt,itemsep=1pt]
  \item a randomly-initialised classifier over the frozen backbone
    (no source training);
  \item a retrieval-augmented (RAG) baseline via DINOv3 feature proximity;
  \item zero-shot transfer of the AIROGS-trained classifier;
  \item a target-supervised oracle trained on labelled on Chaksu data; and
  \item source-free adaptation of the AIROGS classifier on unlabelled Chaksu images via SHOT-style information maximisation.
\end{enumerate}

\paragraph{Principal findings.}
Our results directly contradict the standard narrative of expected failure and necessary adaptation. We discovered that the zero-shot source model
did not fail; it achieved $0.852$ \auroc{} on Indian retinas, only $0.040$ below the
oracle ($0.892$) and far above random-head ($0.487$) and RAG ($0.509$)
baselines. The foundation model already possessed strong inherent cross-ethnic generalisation.

When we subsequently attempted to adapt the model using standard information-maximisation, the performance collapsed to $0.390$ \auroc{}, with $100\%$ sensitivity and $1.4\%$
specificity. The SHOT objective was numerically satisfied by this degenerate
solution, but the patient-level decision rule was completely destroyed. The backbone had
already learned the clinically relevant invariants (optic-disc geometry,
neuroretinal-rim contrast, vessel structure) that survive the
European-to-Indian shift, leaving no useful gap for an unsupervised
objective to close.

We name this failure mode the \textbf{illusion of calibration}: a model
that appears statistically well-behaved in aggregate while losing
individual-level diagnostic discrimination. Feature-space and attention
analyses show that adaptation drifts toward camera-style cues rather than
optic-nerve pathology, because those cues offer an easier optimisation path
for the SFDA loss.

\paragraph{Contributions.} This paper makes three contributions:
\begin{enumerate}[leftmargin=1.2em,topsep=2pt,itemsep=1pt]
  \item \emph{Empirical inversion of the cross-ethnic adaptation narrative.}
    On the largest publicly available Indian glaucoma dataset, DINOv3
    zero-shot transfer reaches near-oracle accuracy while SHOT-style SFDA
    catastrophically degrades the same model.
  \item \emph{Identification and characterisation of the illusion of
    calibration.} A failure mode where aggregate SFDA objectives appear
    satisfied while patient-level discrimination is destroyed, grounded in
    t-SNE and attention-map evidence of camera-confounder takeover.
  \item \emph{A safety-oriented deployment prescription.} On tasks where
    a strong foundation model already generalises, unsupervised adaptation
    must not be trusted without direct labelled validation; aggregate
    calibration metrics cannot detect this failure.
\end{enumerate}
These findings expose a critical vulnerability: deploying unsupervised
test-time adaptation on foundation models for cross-ethnic clinical
applications is currently unsafe.

% ===========================================================
%  2.  RELATED WORK
% ===========================================================
\section{Related Work}
\label{sec:related}

\subsection{Cross-Ethnic Fairness in Ophthalmic AI}

Burlina et al.\ \shortcite{burlina2021} and FairDomain~\cite{fairDomain2024}
established that fundus-based models underperform on darker-skinned patients,
with accuracy gaps of $12.5$ percentage points and \auroc{} disparities of
up to $0.085$ for glaucoma classification, respectively. Wang et al.\ \shortcite{wang2025} 
further demonstrated that domain shifts in diabetic retinopathy screening induce 
significant unfairness, with performance disparities heavily impacting specific ethnic populations.
Coyner et al.\ \shortcite{coyner2023} showed that race-linked signal persists in retinal
representations even after vessel-map abstraction, cautioning that apparently
neutral deep features may still encode subgroup information. 
These findings formed a consensus that models fail across populations, motivating domain adaptation strategies.

Rajesh et al.\ \shortcite{rajesh2025} introduced the Retinal Pigment Score
(RPS), derived from genome-wide association with 20 pigmentation loci. Their
key finding, that RPS overlaps substantially across self-reported ethnic
groups, motivates treating the deployment challenge as a
feature-distribution shift problem rather than a categorical demographic one.
This duality between strong anatomical transferability and latent demographic
confounding is central to the failure mode we identify. Our work directly challenges the prevailing assumption that unsupervised adaptation is always necessary or beneficial when deploying modern foundation models across these phenotypic shifts.

\subsection{Source-Free Domain Adaptation and Test-Time Adaptation}

Source-free domain adaptation (SFDA) and test-time adaptation (TTA)
address the setting where a deployed model must improve on unlabelled
target data without accessing the original training distribution.

\paragraph{Foundational methods.}
SHOT~\cite{liang2020} established the standard SFDA objective: freeze
the classifier head and adapt the feature extractor by minimising
prediction entropy while maximising marginal output diversity. TENT~\cite{wang2021tent} simplified this to updating only batch-normalisation
affine parameters online. Both are appealing in medical settings because
they respect privacy constraints and require no target annotation.

\paragraph{Stability and failure modes.}
Liang et al.\ \shortcite{liang2020} introduce the diversity term in SHOT
precisely because pure entropy minimisation admits degenerate
constant-prediction fixed points; as our results show, however, the
diversity term alone is insufficient under strong label-prior shift. EATA~\cite{niu2022eata}
addresses catastrophic forgetting by filtering high-entropy samples from
gradient updates; SAR~\cite{niu2023sar} extends this to temporally varying
shifts. The shared trajectory of this line of work is explicit recognition
that naive entropy minimisation is unsafe.

\paragraph{Medical imaging.}
Li et al.\ \shortcite{li2025aif} proposed AIF-SFDA using frequency-domain
decomposition for retinal vessel segmentation. The comprehensive survey of
Fang et al.\ \shortcite{fang2024survey} catalogues mixed TTA results across
medical imaging modalities, consistent with our finding that adaptation
outcomes depend sharply on source model strength and the source--target gap.

\subsection{Retinal Foundation Models}

RETFound~\cite{zhou2023} pretrains via masked autoencoders on 1.6 million
unlabelled retinal images, achieving consistent downstream gains across
ophthalmic tasks. DINO-family Vision Transformers~\cite{caron2021} are
particularly attractive for cross-ethnic deployment because their
self-distillation objective promotes globally coherent, shape-sensitive
representations that naturally mitigate pigmentation-driven shift.
DINOv3~\cite{simeoni2025} advances this with rotary positional encoding
and improved patch embedding, outperforming medical-specific foundation
models on multiple 2D imaging tasks despite training exclusively on natural
images. Self-supervised glaucoma detection~\cite{yi2024parcl} confirms
that such networks concentrate attention on the clinically relevant optic
disc and cup, the same behaviour we observe in our zero-shot source model
before adaptation disrupts it.

% ===========================================================
%  3.  METHODOLOGY
% ===========================================================
\section{Methodology}
\label{sec:method}

We describe the five experimental conditions, the auxiliary RAG baseline,
and the shared backbone, training, and evaluation protocol.

\subsection{Problem Formulation}

Let $\mathcal{D}_S = \bigl\{(\mathbf{x}_i, y_i)\bigr\}_{i=1}^{N_S}$ denote the
labelled source domain (\src{}; European tabletop-camera fundus images,
$y \in \{0,1\}$ for non-referable / referable glaucoma) and
$\mathcal{D}_T = \bigl\{\mathbf{x}_j\bigr\}_{j=1}^{N_T}$ the unlabelled
target domain (Chaksu; Indian handheld-camera fundus images). A source
classifier $f_S = h \circ g$ is trained on $\mathcal{D}_S$, where
$g: \mathbb{R}^{3 \times H \times W} \to \mathbb{R}^d$ is a DINOv3
feature extractor and $h: \mathbb{R}^d \to \Delta^2$ is a small
classification head. At deployment, only $f_S$ and unlabelled target
samples are available, and the goal is patient-level glaucoma
discriminability (\auroc{}) on the held-out Chaksu test split.

The compound domain shift comprises retinal pigmentation differences
(\src{} European vs.\ Chaksu South Asian), handheld-versus-tabletop optics,
device-specific artefacts, and a substantial label-prior shift
(\src{} train: $50\%/50\%$; Chaksu test: $15.2\%/84.8\%$).

\subsection{Backbone and Source Training}
\label{ssec:source}

The backbone is the DINOv3 ViT-L/16~\cite{simeoni2025} (303{,}131{,}650
parameters total). We unfreeze the last two transformer blocks together
with the classification head, yielding $25{,}196{,}546$ trainable
parameters ($8.3\%$ of total). The classification head is a two-layer
MLP $d \!\to\! 512 \!\to\! 2$. Source training uses binary cross-entropy
with label smoothing ($\varepsilon = 0.05$), AdamW (lr $= 1\times10^{-4}$,
weight decay $= 0.01$), and a cosine schedule over $50$ epochs with early
stopping on the \src{} validation split. Inputs are resized to
$512 \times 512$ and normalised with ImageNet statistics. Training-time
augmentations are horizontal/vertical flips, $\pm 15^\circ$ rotation,
colour jitter (brightness/contrast $\pm 0.2$, saturation $\pm 0.1$), and
random resized cropping (scale $0.8$--$1.0$).

\subsection{Experimental Conditions}
\label{ssec:conditions}

We evaluate five conditions on the held-out 336-image Chaksu test split
(51 glaucoma, 285 normal).

\paragraph{(1) Random-head baseline.}
The frozen DINOv3 backbone is paired with a randomly initialised two-layer
head (no source training) and evaluated directly on Chaksu. This isolates
how much of the eventual zero-shot performance comes from the trained
classifier head versus the self-supervised backbone alone.

\paragraph{(2) Retrieval-augmented (RAG) baseline.}
For each Chaksu test image, we retrieve its $k$ nearest neighbours from a
FAISS L2 index built over DINOv3 CLS-token embeddings of the \src{} test
set, and predict by majority vote. We grid-search $k \in \{5, 10, 20, 50\}$
and three aggregation methods (majority vote, weighted vote, mean
similarity). This tests whether cross-ethnic label transfer is viable
through feature-space proximity alone, without any classifier.

\paragraph{(3) Zero-shot source transfer.}
$f_S$, trained on \src{} only, is applied directly to Chaksu with no
target-side updates. This is the headline condition: it tests the
prevailing claim that European-trained models fail on Indian retinas.

\paragraph{(4) Target-supervised oracle.}
A classifier with the same architecture is trained under the same protocol
as \S\ref{ssec:source}, but on labelled Chaksu training data (the official
1{,}009-image Chaksu training split). The oracle establishes an upper-bound
reference and quantifies the remaining gap that any unsupervised adaptation
would need to close.

\paragraph{(5) Source-free adaptation (SHOT-style IM).}
\label{par:shot}
Starting from $f_S$, we adapt the feature extractor $g$ on the unlabelled
Chaksu training images using the SHOT~\cite{liang2020} information-maximisation
objective:
\begin{equation}
  \mathcal{L}_\mathrm{IM} \;=\;
    \underbrace{-\sum_k \hat{p}_k \log \hat{p}_k}_{\mathcal{L}_\mathrm{ent}}
    \;-\; \lambda \underbrace{\left(-\sum_k \bar{p}_k \log \bar{p}_k\right)}_{\mathcal{L}_\mathrm{div}},
  \label{eq:shot}
\end{equation}
where $\hat{p}_k$ is the per-sample softmax output and
$\bar{p}_k = \tfrac{1}{B}\sum_i \hat{p}_k^{(i)}$ is the batch-mean
prediction. We additionally apply a prevalence-matching regulariser:
\begin{equation}
  \mathcal{L}_\mathrm{prev} \;=\;
    D_\mathrm{KL}\!\left(\pi \,\big\|\, \bar{p}_\mathrm{glaucoma}\right),
  \label{eq:prev}
\end{equation}
giving a total objective $\mathcal{L} = \mathcal{L}_\mathrm{IM} +
\mu\,\mathcal{L}_\mathrm{prev}$ with $\mu = 0.1$, anchoring the
batch-mean glaucoma prediction to the estimated target prior
$\pi \approx 0.152$. The classification head $h$ is frozen; the feature
extractor $g$ is adapted for $10$ epochs with Adam
(lr $= 1\times10^{-4}$, batch size $32$).

% ===========================================================
%  4.  EVALUATION
% ===========================================================
\section{Evaluation}
\label{sec:eval}

\subsection{Datasets}

\paragraph{\src{} (Source Domain).}
We use the balanced EyePACS-AIROGS-light-V2 subset~\cite{devente2024}:
$9{,}540$ preprocessed $512\times512$ fundus photographs split evenly
between referable and non-referable glaucoma, with 4{,}000/385/385
images per class for train/validation/test. The images originate from
Dutch population-based screening and were standardised to remove large
background regions before training.

\paragraph{Chaksu (Target Domain).}
Chaksu~\cite{kumar2023} is the largest publicly available
Indian-ethnicity glaucoma fundus database: $1{,}345$ images from $674$
patients annotated by fellowship-trained glaucoma specialists. The public
release provides 1{,}009 training images and 336 held-out test images
(51 glaucoma, 285 normal), acquired with three handheld devices (Remidio, Forus, Bosch) with substantial variation in colour profile,
sharpness, field-of-view, and visible device artefacts.

\subsection{Evaluation Metrics}

We report the following on the held-out 336-image Chaksu test split:
\begin{itemize}[leftmargin=1.2em,topsep=2pt,itemsep=1pt]
  \item \textbf{\auroc{}} \textendash{} primary metric; threshold-independent.
  \item \textbf{Sensitivity (\se{})} and \textbf{Specificity (\spec{})}
    \textendash{} at the Youden-optimal threshold
    $\hat{\theta} = \arg\max_\theta (\mathrm{Se}(\theta) + \mathrm{Sp}(\theta) - 1)$.
  \item \textbf{F1 Score} \textendash{} harmonic mean of precision and recall
    at $\hat{\theta}$, important under the $51/285$ class imbalance.
  \item \textbf{Sens@95\%Spec} \textendash{} sensitivity at a fixed
    high-specificity operating point relevant to first-pass screening.
\end{itemize}
Accuracy is not reported as a primary metric: a constant ``normal''
predictor achieves $84.8\%$ accuracy with zero clinical utility.

\subsection{Implementation Details}

All experiments were run on an RTX 5090 with 32GB VRAM. The
DINOv3 ViT-L/16 backbone is initialised from the official
release~\cite{simeoni2025}. Chaksu images arrive at
$1920\!\times\!1440$ to $2448\!\times\!3264$ resolution; aggressive
preprocessing (circular fundus extraction, black-border removal) precedes
resizing to $512\times512$. Full hyperparameter settings are provided in
Appendix~\ref{app:hyperparams}.

\subsection{Main Results}

Table~\ref{tab:main} reports all five conditions on the held-out Chaksu
test split, plus the \src{} source-domain sanity check ($\auroc{} = 0.980$).

\begin{table*}[t]
\centering
\small
\caption{Main results on the 336-image Chaksu held-out test split (51
  glaucoma, 285 normal). Sensitivity, specificity, F1, and Sens@95\%Spec
  are reported at the Youden-optimal threshold (Sens@95\%Spec at the
  fixed-specificity threshold). The AIROGS source-domain row is a
  sanity check confirming that the source classifier is well trained.
  Best result among source-free conditions in bold.}
\label{tab:main}
\begin{tabular}{lccccccc}
\toprule
\textbf{Condition} & \textbf{Source-free?} & \textbf{\auroc{}} & \textbf{\se{} (\%)} & \textbf{\spec{} (\%)} & \textbf{Prec.\ (\%)} & \textbf{F1} & \textbf{Sens@95\%Sp} \\
\midrule
\multicolumn{8}{l}{\emph{\src{} sanity check}} \\
AIROGS\,$\to$\,AIROGS  & --- & 0.980 & --- & --- & --- & --- & --- \\
\midrule
\multicolumn{8}{l}{\emph{Cross-ethnic deployment on Chaksu}} \\
Random Head (no source training)        & Yes & 0.487 & 19.6 & 90.5 & 27.0 & 0.227 & 11.8 \\
Retrieval (RAG, $k{=}5$, majority)      & Yes & 0.509 & 25.0 & 74.0 & --- & 0.231 & --- \\
Zero-Shot Source Transfer               & Yes & \textbf{0.852} & 76.5 & \textbf{83.9} & 45.9 & \textbf{0.574} & \textbf{49.0} \\
SHOT-IM Source-Free Adaptation          & Yes & 0.390 & 100.0 & 1.4 & 15.4 & 0.266 & 2.0 \\
\midrule
\textit{Target-Supervised Oracle}       & \textit{No} & \textit{0.892} & \textit{92.2} & \textit{78.2} & \textit{43.1} & \textit{0.588} & \textit{45.1} \\
\bottomrule
\end{tabular}
\end{table*}

\paragraph{The source model does not fail on Indian retinas.}
The zero-shot classifier achieves $0.852$ \auroc{} on Chaksu, only $0.040$
below the target-supervised oracle. Given that $95\%$ bootstrap confidence
intervals on \auroc{} at this sample size are approximately $\pm 0.04$, this
gap should be read as an upper-bound estimate. The zero-shot model delivers
a clinically usable operating point ($76.5\%$ sensitivity, $83.9\%$
specificity, F1 $= 0.574$, Sens@95\%Spec $= 49.0\%$) without any access
to target labels during training. Notably, the oracle achieves \emph{lower}
specificity ($78.2\%$) than the zero-shot model ($83.9\%$) despite access
to $1{,}009$ labelled Chaksu images, suggesting it has partially overfit
to camera-specific cues (the same confounders that later cause SFDA
collapse), while the zero-shot model retains a device-agnostic
disc-geometry boundary.

\paragraph{Transfer requires the source classifier, not just the backbone.}
The random-head baseline achieves only $0.487$ \auroc{}, essentially chance.
The RAG baseline, assigning labels by majority vote over $k$ AIROGS nearest
neighbours in DINOv3 feature space, performs equally poorly ($0.509$ at
$k{=}5$; full grid over $k \in \{5,10,20,50\}$ spans $0.457$--$0.509$).
Cross-ethnic transfer is not a free property of the backbone or of
nearest-neighbour similarity: it requires the source-supervised classifier
head that has learned to map DINOv3 features to the glaucoma decision
boundary.

\paragraph{Unsupervised adaptation collapses catastrophically.}
SHOT-style information-maximisation adaptation degrades \auroc{} to
$0.390$, which is worse than chance and $0.462$ below the zero-shot
baseline it was meant to improve. The adapted model predicts glaucoma for $98.6\%$ of
test images ($100\%$ sensitivity, $1.4\%$ specificity), clinically
equivalent to a constant glaucoma predictor. Crucially, the SHOT objective
is \emph{numerically satisfied}: batch-mean predictions match the target
prevalence and per-sample entropies are low. That is precisely what makes
this failure dangerous in deployment.

\subsection{The Illusion of Calibration}
\label{ssec:illusion}

\paragraph{Mechanism.}
The SHOT-IM objective in Eq.~\ref{eq:shot} is \emph{satisfiable by a
degenerate decision rule}. The optimiser finds the path of least
resistance: exploit the most salient target-domain features (device
lighting artefacts, melanin-driven colour, watermarks) to satisfy the
prevalence constraint, rather than learning glaucoma pathology. The result
is a model that produces the correct marginal class distribution while
assigning labels to entirely wrong samples. Formally, this satisfies
$D_{\mathrm{KL}}(\bar{p}\,\|\,\pi)\approx 0$ while simultaneously
maximising the per-sample misclassification rate---possible because the
information-maximisation objective provides no pairwise or structural
constraint linking feature similarity to label assignment. The label-prior
shift (\src{} train: $50\%/50\%$; Chaksu test: $15.2\%/84.8\%$) further
biases adaptation toward confident, pathology-agnostic fixed points.

\paragraph{Feature-space evidence.}
Figure~\ref{fig:tsne} compares t-SNE projections of penultimate-layer
features from the zero-shot and SHOT-adapted models. The zero-shot
embedding retains visible class structure consistent with its $0.852$
\auroc{}. After adaptation, glaucomatous and normal cases become
substantially less separable, consistent with the $0.390$ \auroc{}.

\begin{figure*}[t]
\centering
\begin{subfigure}[t]{0.46\textwidth}
  \centering
  \includegraphics[width=\linewidth]{Evaluation-Run-7/tsne_source.png}
  \caption{Zero-shot source model.}
\end{subfigure}\hfill
\begin{subfigure}[t]{0.46\textwidth}
  \centering
  \includegraphics[width=\linewidth]{Evaluation-Run-7/tsne_adapted.png}
  \caption{SHOT-adapted model.}
\end{subfigure}
\caption{t-SNE projections of penultimate-layer DINOv3 features on the
  Chaksu test split. Class-discriminative structure visible in the
  zero-shot embedding (left) is largely lost after SHOT-IM adaptation
  (right).}
\label{fig:tsne}
\end{figure*}

\paragraph{Attention evidence.}
Figure~\ref{fig:attention} shows DINOv3 \texttt{[CLS]}-token attention
overlays on representative Chaksu images. The zero-shot model concentrates
on the optic-disc and cup, the clinically relevant anatomy for glaucoma
assessment, consistent with Yi et al.\ \shortcite{yi2024parcl}.
The adapted model's attention diffuses outward to peripheral retina,
vessels, and image edges that covary systematically with camera identity
(illumination falloff, field-of-view, watermark borders) but carry no
direct glaucomatous signal.

\begin{figure*}[t]
\centering
\includegraphics[width=0.98\textwidth]{Evaluation-Run-7/attention_maps_comparison.png}
\caption{DINOv3 \texttt{[CLS]} attention overlays on representative
  Chaksu images. Top: original. Middle: zero-shot source model
  (attention centred on the optic disc/cup). Bottom: SHOT-adapted model
  (diffuse attention drifting away from the clinically relevant
  region).}
\label{fig:attention}
\end{figure*}

\paragraph{Why the foundation model already generalises.}
DINOv3's self-distillation objective compels the model to extract
geometric invariants of the optic nerve head (disc/cup depth, rim width,
vessel curvature) rather than brittle colour and texture cues that
narrower architectures overfit to~\cite{caron2021,simeoni2025}. Source
supervision on \src{} then aligns these shape-sensitive features with the
glaucoma decision boundary. Because optic-disc geometry is preserved
across European and Indian retinas irrespective of pigmentation, the
zero-shot decision rule transfers. Unsupervised adaptation, by contrast,
has no signal that prefers disc geometry; it moves features toward whatever
reduces the SFDA loss fastest, and camera-style cues do that far more
easily than subtle structural pathology.

\subsection{Training Dynamics}

Figure~\ref{fig:losscurves} shows training loss curves for the source,
oracle, and SHOT-IM conditions. The source and oracle curves follow the
expected smooth supervised descent. The SHOT-IM curve descends rapidly to
a low value; the adaptation objective \emph{is} minimised. That is
precisely what makes the resulting $0.390$ \auroc{} diagnostically
dangerous: the optimiser does its job, but on the wrong objective.

\begin{figure*}[t]
\centering
\begin{subfigure}[t]{0.32\textwidth}
  \centering
  \includegraphics[width=\linewidth]{Evaluation-Run-7/Source_AIROGS/loss_curve_source.png}
  \caption{Source (AIROGS).}
\end{subfigure}\hfill
\begin{subfigure}[t]{0.32\textwidth}
  \centering
  \includegraphics[width=\linewidth]{Evaluation-Run-7/Oracle_Chakshu/loss_curve_oracle.png}
  \caption{Oracle (Chaksu).}
\end{subfigure}\hfill
\begin{subfigure}[t]{0.32\textwidth}
  \centering
  \includegraphics[width=\linewidth]{Evaluation-Run-7/NetraAdapt/loss_curve.png}
  \caption{SHOT-IM adapt.}
\end{subfigure}
\caption{Training loss curves. Source and oracle conditions follow the
  expected supervised descent; SHOT-IM adaptation descends to a low
  objective value while test \auroc{} simultaneously collapses
  (Table~\ref{tab:main}).}
\label{fig:losscurves}
\end{figure*}

% ===========================================================
%  5.  CONCLUSION
% ===========================================================
\section{Conclusion}
\label{sec:conclusion}

We directly test two widely held assumptions about cross-ethnic deployment of
retinal AI: that European-trained models significantly degrade on Indian retinas, and that
unsupervised source-free adaptation is the necessary fix. Neither holds.

A DINOv3 ViT-L/16 source-trained on \src{} reaches $0.852$ \auroc{} on
the held-out Chaksu test split, only $0.040$ below a target-supervised
oracle, without ever seeing a target image during training. Random-head
($0.487$) and retrieval-based ($0.509$) baselines remain near chance,
confirming this is not a free backbone property but the result of a
source-trained classifier that has learned the glaucoma decision boundary.
The compound phenotype-and-device shift between European and Indian
captures is real, but the source model has already learned the
representations (optic-disc geometry, rim contrast, vessel structure)
that survive it.

Applying SHOT-style information-maximisation adaptation to this well-generalising model causes it to collapse to $0.390$ \auroc{} with $100\%$ sensitivity and $1.4\%$
specificity. The adaptation objective is successfully minimised, yet the patient-level
decision rule is destroyed. We name this failure mode the \emph{illusion
of calibration}, and ground its mechanism in feature-space and attention
evidence: lacking any signal that prefers optic-disc anatomy over
camera-style cues, the adapted model drifts toward the easier path.

We draw three principal conclusions:
\begin{enumerate}[leftmargin=1.2em,topsep=2pt,itemsep=1pt]
  \item \textbf{Foundation models inherently generalise across ethnic shifts.} A foundation model trained on $9{,}540$ European images
    provides more clinical utility zero-shot than SFDA on $1{,}009$
    unlabelled Indian images. The cross-ethnic gap is small enough that
    unsupervised adaptation is unnecessary and harmful.
  \item \textbf{Unsupervised TTA is currently unsafe for cross-ethnic
    medical deployment.} Algorithms that align unlabelled feature
    distributions cannot distinguish a shift in camera lighting from a
    shift in disease severity. The illusion of calibration satisfies
    aggregate statistics while destroying individual discrimination;
    only a labelled held-out evaluation can detect this failure.
  \item \textbf{Feature retrieval cannot bridge ethnic gaps.} The RAG
    baseline ($0.509$ \auroc{}) confirms that general visual similarity
    in DINOv3 feature space does not capture pathological similarity
    across phenotypes; a supervised decision boundary is essential.
\end{enumerate}
Future work should move beyond pure entropy-based objectives toward
structurally constrained adaptation (e.g., optic-disc-anchored feature
alignment) or few-shot supervised techniques to safely extend foundation
models across demographic boundaries in ophthalmology.

% ===========================================================
%  6.  LIMITATIONS AND FUTURE WORK
% ===========================================================
\section{Limitations and Future Work}
\label{sec:futurelimitations}

\paragraph{Scope of the finding.}
Our study uses a single source--target pair (\src{}$\to$Chaksu); whether
the same inversion holds for other pairings (e.g., AIROGS$\to$African or
East Asian cohorts) requires further study. Where the source distribution
is narrower or the cross-ethnic gap is larger, the failure premise may
actually hold and adaptation may be more useful. Results should not be
generalised beyond the setting tested.

\paragraph{Single SFDA family.}
We test the canonical SHOT information-maximisation objective; other
families (TENT~\cite{wang2021tent}, EATA~\cite{niu2022eata},
SAR~\cite{niu2023sar}, AIF-SFDA~\cite{li2025aif}) may fail less
catastrophically, but the core issue (that the unlabelled-target
objective is satisfiable by camera-style confounders) is not specific
to SHOT.

\paragraph{Mechanistic evidence and replication.}
Our t-SNE and attention analyses are visual and qualitative; controlled
counterfactual studies (e.g., camera-artefact ablation) would provide
causal evidence. Each condition is run once; SHOT-IM collapse magnitude
may vary across seeds and batch orderings.

\paragraph{Directions forward.}
The most practically valuable follow-up is few-shot adaptation: even
$5$--$50$ labelled target images per class would anchor the decision
boundary to actual pathology and structurally prevent a constant-predictor
collapse. Stratifying zero-shot performance by the Retinal Pigment
Score~\cite{rajesh2025} rather than by self-reported ethnicity would
provide a more biologically grounded characterisation of the transfer
behaviour we observe.

% ===========================================================
%  ACKNOWLEDGMENTS
% ===========================================================
\section*{Acknowledgments}

We thank the Chaksu dataset creators at the Indian Institute of Science, Bangalore, and the AIROGS challenge organisers for making their datasets publicly available, without which this study would not have been possible.

% ===========================================================
%  BIBLIOGRAPHY
% ===========================================================
\begin{thebibliography}{19}

\bibitem[Burlina et~al.(2021)]{burlina2021}
Philip Burlina, Neil Joshi, William Paul, Kapil D. Pacheco, and Neil M. Bressler.
\newblock 2021.
\newblock Addressing artificial intelligence bias in retinal diagnostics.
\newblock \textit{Translational Vision Science \& Technology}, 10(2):13.

\bibitem[Caron et~al.(2021)]{caron2021}
Mathilde Caron, Hugo Touvron, Ishan Misra, Hervé Jégou, Julien Mairal,
  Piotr Bojanowski, and Armand Joulin.
\newblock 2021.
\newblock Emerging properties in self-supervised vision transformers.
\newblock In \textit{Proceedings of the IEEE/CVF International Conference on
  Computer Vision (ICCV)}, pages 9650--9660.

\bibitem[Coyner et~al.(2023)]{coyner2023}
Aaron S. Coyner, Praveer Singh, James M. Brown, Susan Ostmo, R.V. Paul Chan,
  Michael F. Chiang, Jayashree Kalpathy-Cramer, J. Peter Campbell, and the
  Imaging and Informatics in Retinopathy of Prematurity Consortium.
\newblock 2023.
\newblock Association of biomarker-based artificial intelligence with risk of
  racial bias in retinal images.
\newblock \textit{JAMA Ophthalmology}, 141(6):543--552.

\bibitem[de~Vente et~al.(2024)]{devente2024}
Coen de~Vente, Koenraad A. Vermeer, Nicolas Jaccard, et~al.
\newblock 2024.
\newblock {AIROGS}: Artificial intelligence for robust glaucoma screening
  challenge.
\newblock \textit{IEEE Transactions on Medical Imaging}, 43(1):542--557.

\bibitem[Fang et~al.(2024)]{fang2024survey}
Yuqi Fang, Haofei Yao, Hao Zhao, Zhilin Zhang, Zikang Ma, Yan Sun, and
  Xiu Li.
\newblock 2024.
\newblock A comprehensive survey on source-free domain adaptation.
\newblock \textit{IEEE Transactions on Pattern Analysis and Machine
  Intelligence (TPAMI)}, 46(8):5743--5762.

\bibitem[Kumar et~al.(2023)]{kumar2023}
J.R. Harish Kumar, Chandra Sekhar Seelamantula, J.H. Gagan, et~al.
\newblock 2023.
\newblock Chaksu: A glaucoma specific fundus image database.
\newblock \textit{Scientific Data}, 10:70.

\bibitem[Li et~al.(2025)]{li2025aif}
Yu Li, Hao Chen, Zhitong Wang, and Yazhou Zhang.
\newblock 2025.
\newblock {AIF-SFDA}: Autonomous information filter-based source-free domain
  adaptation for medical image segmentation.
\newblock In \textit{Proceedings of the AAAI Conference on Artificial
  Intelligence (AAAI)}.

\bibitem[Liang et~al.(2020)]{liang2020}
Jian Liang, Dapeng Hu, and Jiashi Feng.
\newblock 2020.
\newblock Do we really need to access the source data? source hypothesis
  transfer for unsupervised domain adaptation.
\newblock In \textit{Proceedings of the 37th International Conference on
  Machine Learning (ICML)}, pages 6028--6039.

\bibitem[Niu et~al.(2022)]{niu2022eata}
Shuaicheng Niu, Jiaxiang Wu, Yifan Zhang, Yaofo Chen, Shijian Zheng,
  Peilin Zhao, and Mingkui Tan.
\newblock 2022.
\newblock Efficient test-time model adaptation without forgetting.
\newblock In \textit{Proceedings of the 39th International Conference on
  Machine Learning (ICML)}, pages 16888--16905.

\bibitem[Niu et~al.(2023)]{niu2023sar}
Shuaicheng Niu, Jiaxiang Wu, Yifan Zhang, Zhiquan Wen, Zhaozheng Chen,
  Peilin Zhao, and Mingkui Tan.
\newblock 2023.
\newblock Towards stable test-time adaptation in dynamic wild world.
\newblock In \textit{International Conference on Learning Representations
  (ICLR)}.

\bibitem[Rajesh et~al.(2025)]{rajesh2025}
Anand E. Rajesh, Abraham Olvera-Barrios, Alasdair N. Warwick, et~al.
\newblock 2025.
\newblock Machine learning derived retinal pigment score from ophthalmic
  imaging shows ethnicity is not biology.
\newblock \textit{Nature Communications}, 16:60.

\bibitem[RETFound-Green(2025)]{retfoundgreen2025}
Yukun Zhou, Siegfried K. Wagner, Mark A. Chia, et~al.
\newblock 2025.
\newblock {RETFound-Green}: Efficient retinal foundation model using {DINOv2}
  pretraining with token reconstruction.
\newblock \textit{Nature Communications}, 16:1242.

\bibitem[Simeoni et~al.(2025)]{simeoni2025}
Oriane Simeoni, Huy V. Vo, Maximilian Seitzer, et~al.
\newblock 2025.
\newblock {DINOv3}.
\newblock \textit{arXiv preprint arXiv:2508.10104}.

\bibitem[Tham et~al.(2014)]{tham2014}
Yih-Chung Tham, Xiang Li, Tien Y. Wong, Harry A. Quigley, Tin Aung, and
  Ching-Yu Cheng.
\newblock 2014.
\newblock Global prevalence of glaucoma and projections of glaucoma burden
  through 2040: A systematic review and meta-analysis.
\newblock \textit{Ophthalmology}, 121(11):2081--2090.

\bibitem[Tian et~al.(2024)]{fairDomain2024}
Yuan Tian, Bing Jiang, Wei Lu, and Kai Ma.
\newblock 2024.
\newblock {FairDomain}: Achieving fairness in cross-domain medical image
  segmentation and classification.
\newblock In \textit{Proceedings of the European Conference on Computer Vision
  (ECCV)}.

\bibitem[Wang et~al.(2021)]{wang2021tent}
Dequan Wang, Evan Shelhamer, Shaoteng Liu, Bruno Olshausen, and Trevor Darrell.
\newblock 2021.
\newblock Tent: Fully test-time adaptation by entropy minimization.
\newblock In \textit{International Conference on Learning Representations
  (ICLR)}.

\bibitem[Wang et~al.(2025)]{wang2025}
Yvette Wang, Zhengquan Luo, Xiaotong Han, Mingguang He, and Chi Liu.
\newblock 2025.
\newblock Domain shift-induced unfairness in {AI} for diabetic retinopathy screening.
\newblock \textit{Investigative Ophthalmology \& Visual Science}, 66(8):3891.

\bibitem[Yi et~al.(2024)]{yi2024parcl}
Xuelin Yi, Long Chen, and Yang Zhang.
\newblock 2024.
\newblock {PaRCL}: Pathology-aware representation contrastive learning for
  glaucoma classification.
\newblock In \textit{Proceedings of the IEEE International Conference on
  Bioinformatics and Biomedicine (BIBM)}.

\bibitem[Zhou et~al.(2023)]{zhou2023}
Yukun Zhou, Mark A. Chia, Siegfried K. Wagner, et~al.
\newblock 2023.
\newblock A foundation model for generalizable disease detection from retinal
  images.
\newblock \textit{Nature}, 622:156--163.

\end{thebibliography}

% ===========================================================
%  APPENDIX
% ===========================================================
\clearpage
\appendix

\section{Hyperparameter Settings}
\label{app:hyperparams}

Table~\ref{tab:hyper} reports the hyperparameter settings used for the
three trained conditions (source, oracle, SHOT-IM adaptation),
matching the configuration logged for the runs that produced
Table~\ref{tab:main}.

\begin{table}[ht]
\centering
\small
\caption{Hyperparameter settings for the trained conditions.}
\label{tab:hyper}
\resizebox{\columnwidth}{!}{%
\begin{tabular}{lll}
\toprule
\textbf{Condition} & \textbf{Hyperparameter} & \textbf{Value} \\
\midrule
\multirow{8}{*}{Source training (\src{})}
  & Backbone & DINOv3 ViT-L/16 \\
  & Trainable params & Last 2 ViT blocks $+$ head \\
  & Trainable / total & 25.2\,M / 303.1\,M ($8.3\%$) \\
  & Loss & BCE w/ label smoothing $\varepsilon = 0.05$ \\
  & Optimiser & AdamW \\
  & Learning rate & $1\times10^{-4}$ \\
  & Weight decay & $0.01$ \\
  & Schedule / epochs & cosine / $50$ (early stop) \\
\midrule
\multirow{4}{*}{Oracle (Chaksu)}
  & Architecture & same as source \\
  & Train/test split & official $1{,}009/336$ \\
  & Optimiser / lr & AdamW / $1\times10^{-4}$ \\
  & Schedule / epochs & cosine / $50$ (early stop) \\
\midrule
\multirow{6}{*}{SHOT-IM adaptation}
  & Frozen / adapted & head / feature extractor \\
  & Optimiser & Adam \\
  & Learning rate & $1\times10^{-4}$ \\
  & Weight decay & $0$ \\
  & Batch size & $32$ \\
  & Epochs & $10$ \\
\midrule
\multirow{4}{*}{Shared}
  & Input resolution & $512 \times 512$ \\
  & Normalisation & ImageNet mean/std \\
  & Train aug.\ & H/V flip, $\pm 15^\circ$, jitter, RRC \\
  & Hardware & 1$\times$ NVIDIA A100 40\,GB \\
\bottomrule
\end{tabular}%
}
\end{table}

\section{RAG Baseline: Full Results}
\label{app:rag}

The retrieval-augmented baseline indexes DINOv3 \texttt{[CLS]}-token
embeddings of the AIROGS test set ($n = 770$) with a FAISS
\texttt{IndexFlatL2}. For each Chaksu test image we retrieve $k$
nearest neighbours and predict via three aggregation methods: majority
vote, similarity-weighted vote, and mean similarity. We grid-search
$k \in \{5, 10, 20, 50\}$, giving $12$ configurations in total. The
best configuration ($k = 5$, majority vote) achieves $0.509$
\auroc{}; \auroc{} across the full grid spans $0.457$--$0.509$,
confirming that the failure of nearest-neighbour transfer is not
sensitive to $k$ or aggregation choice.

Table~\ref{tab:rag_full} reports all 12 configurations.
Figure~\ref{fig:rag_heatmap} shows the AUROC and sensitivity--specificity
landscape as a heatmap. Figure~\ref{fig:rag_best} shows the ROC curve and
confusion matrix for the best configuration ($k{=}5$, majority vote).

\begin{table}[H]
\centering
\small
\setlength{\abovecaptionskip}{4pt}
\setlength{\belowcaptionskip}{4pt}
\caption{Full RAG grid-search results. The RAG pipeline was evaluated
         on an earlier Chaksu preprocessing run yielding 302 images
         (233 normal, 69 glaucoma), compared to the 336-image split
         (285/51) used for the DNN conditions in Table~\ref{tab:main}.
         All AUROC values are near 0.5, confirming retrieval-based
         transfer fails regardless of hyperparameter choice.}
\label{tab:rag_full}
\resizebox{\columnwidth}{!}{%
\begin{tabular}{llcccc}
\toprule
$k$ & Aggregation & \auroc{} & Se & Sp & F1 \\
\midrule
 5 & Majority vote    & 0.509 & 0.246 & 0.738 & 0.231 \\
 5 & Weighted vote    & 0.505 & 0.246 & 0.738 & 0.231 \\
 5 & Mean prob.       & 0.509 & 0.246 & 0.738 & 0.231 \\
10 & Majority vote    & 0.501 & 0.304 & 0.674 & 0.253 \\
10 & Weighted vote    & 0.501 & 0.246 & 0.755 & 0.238 \\
10 & Mean prob.       & 0.501 & 0.304 & 0.674 & 0.253 \\
20 & Majority vote    & 0.462 & 0.275 & 0.747 & 0.259 \\
20 & Weighted vote    & 0.465 & 0.232 & 0.794 & 0.241 \\
20 & Mean prob.       & 0.462 & 0.275 & 0.747 & 0.259 \\
50 & Majority vote    & 0.458 & 0.217 & 0.773 & 0.219 \\
50 & Weighted vote    & 0.458 & 0.203 & 0.803 & 0.217 \\
50 & Mean prob.       & 0.458 & 0.217 & 0.773 & 0.219 \\
\bottomrule
\end{tabular}%
}
\end{table}

\begin{figure}[H]
  \centering
  \setlength{\abovecaptionskip}{4pt}
  \begin{subfigure}[t]{0.48\columnwidth}
    \centering
    \includegraphics[width=\linewidth]{auroc_heatmap.pdf}
    \caption{AUROC across the $k \times$ method grid.
             All values cluster near 0.5 (random chance),
             demonstrating insensitivity to hyperparameters.}
  \end{subfigure}\hfill
  \begin{subfigure}[t]{0.48\columnwidth}
    \centering
    \includegraphics[width=\linewidth]{sensitivity_specificity.pdf}
    \caption{Sensitivity--specificity trade-off across configurations.
             Larger $k$ inflates specificity at the cost of sensitivity;
             no configuration achieves both above 0.5.}
  \end{subfigure}
  \caption{RAG grid-search landscape. The heatmap and trade-off plot
           confirm that the retrieval approach does not have a useful
           operating regime in this cross-ethnic setting.}
  \label{fig:rag_heatmap}
\end{figure}

\begin{figure}[H]
  \centering
  \setlength{\abovecaptionskip}{4pt}
  \begin{subfigure}[t]{0.48\columnwidth}
    \centering
    \includegraphics[width=\linewidth]{roc_curve.pdf}
    \caption{ROC curve for the best RAG configuration ($k{=}5$, majority
             vote). The curve hugs the diagonal, confirming performance
             indistinguishable from a random classifier.}
  \end{subfigure}\hfill
  \begin{subfigure}[t]{0.48\columnwidth}
    \centering
    \includegraphics[width=\linewidth]{confusion_matrix.pdf}
    \caption{Confusion matrix. The system correctly identifies normals
             (TN\,=\,172) but misses 75\% of glaucoma cases
             (FN\,=\,52 out of 69), making it clinically unusable.}
  \end{subfigure}
  \caption{Best-configuration RAG results ($k{=}5$, majority vote,
           AUROC\,$= 0.509$). Even the best configuration misses
           three in four glaucoma patients.}
  \label{fig:rag_best}
\end{figure}

\end{document}